\section{The A* algorithm}\label{sec:astar}

A*~\cite{hart1968formal} keeps the cost so far of Dijkstra's algorithm and adds
the estimate of greedy best-first search (Figure~\ref{fig:combine}).

\begin{figure}[htbp]
  \centering
  \tikzset{
  card/.style  = {draw=#1, fill=#1!9, line width=0.9pt, rounded corners=3pt,
                  text width=4.8cm, align=center, font=\small, inner sep=6pt},
  cflow/.style = {-{Stealth[length=2.4mm]}, line width=1pt, draw=inkgray},
  cflab/.style = {font=\small, text=inkgray, fill=white, inner sep=2pt},
}
\begin{tikzpicture}[x=1cm, y=1cm]
  \node[card=algdijkstra] (dj) at (-3.4,1.7)
    {\textcolor{algdijkstra}{\textbf{Dijkstra's algorithm}}\\[2pt]
     ranks by $g$, the cost so far\\[3pt]
     \textcolor{algastar}{$\checkmark$} shortest path\\
     \textcolor{cbRed}{$\times$} expands all eight cities};
  \node[card=alggreedy] (gr) at (3.4,1.7)
    {\textcolor{alggreedy}{\textbf{Greedy best-first search}}\\[2pt]
     ranks by $h$, the estimate of what is left\\[3pt]
     \textcolor{algastar}{$\checkmark$} expands four cities\\
     \textcolor{cbRed}{$\times$} path cost 10, optimum 8};
  \node[card=algastar, line width=1.3pt, text width=5cm] (as) at (0,-1.6)
    {\textcolor{algastar}{\textbf{A* search}}\\[2pt]
     ranks by $f = g + h$\\[3pt]
     \textcolor{algastar}{$\checkmark$} shortest path (cost 8)\\
     \textcolor{algastar}{$\checkmark$} expands five cities};
  \draw[cflow] (dj.south) -- node[cflab] {keep $g$} ($(as.north)+(-1.2,0)$);
  \draw[cflow] (gr.south) -- node[cflab] {add $h$} ($(as.north)+(1.2,0)$);
\end{tikzpicture}

  \caption{A* takes the cost so far $g$ from Dijkstra's algorithm and the
    estimate $h$ from greedy best-first search. The numbers are those of the
    running example.}\label{fig:combine}
\end{figure}

\subsection{Combining the two keys}\label{sec:fgh}

A* ranks every vertex $v$ of $Q$ by
\begin{equation}\label{eq:f}
  f(v) = g(v) + h(v),
\end{equation}
the estimated cost of the cheapest complete route through $v$
(Figure~\ref{fig:fgh}). A* thus prefers the vertex on the most promising
complete route, rather than the vertex nearest the source or the one that looks
nearest the destination (Table~\ref{tab:three-keys}).

\begin{figure}[htbp]
  \centering
  \begin{tikzpicture}[x=1cm, y=1cm]
    \node[storystart] (S) at (0,0) {$s$};
    \node[storynode]  (V) at (4,0) {$v$};
    \node[storygoal]  (T) at (10,0) {$t$};
    \draw[algastar, line width=1.4pt] (S) -- (V);
    \draw[alggreedy, line width=1.4pt, dash pattern=on 5pt off 3pt] (V) -- (T);
    \draw[algastar, {Stealth[length=2mm]}-{Stealth[length=2mm]}]
      (0,-0.6) -- node[below=2pt, font=\small, align=center] {$g(v)$\\cost of
      the best known path} (3.95,-0.6);
    \draw[alggreedy, {Stealth[length=2mm]}-{Stealth[length=2mm]}]
      (4.05,-0.6) -- node[below=2pt, font=\small, align=center] {$h(v)$\\
      estimate of the remaining cost} (10,-0.6);
    \draw[rbluedk, {Stealth[length=2mm]}-{Stealth[length=2mm]}]
      (0,0.6) -- node[above=2pt, font=\small] {$f(v) = g(v) + h(v)$:
      estimated cost of the whole path through $v$} (10,0.6);
  \end{tikzpicture}
  \caption{The key of A*. The first part is known exactly; the second part is
    an estimate.}\label{fig:fgh}
\end{figure}

\begin{table}[htbp]
  \centering
  \caption{The key of each algorithm.}\label{tab:three-keys}
  \begin{tabular}{@{}lll@{}}
    \toprule
    Algorithm & Key of $v$ & What the key measures\\
    \midrule
    Dijkstra's algorithm     & $g(v)$        & the cost already paid\\
    Greedy best-first search & $h(v)$        & the estimated cost still to pay\\
    A* search                & $g(v) + h(v)$ & the estimated cost of the whole path\\
    \bottomrule
  \end{tabular}
\end{table}

\subsection{A geometric view of the key}\label{sec:geometry}

On an open plane, let $g$ be the distance from $s$ and $h$ the distance to $t$.
Then $g$ and $h$ are cones around $s$ and $t$, and their sum $f$ is a trough
whose floor is the segment from $s$ to $t$ (Figure~\ref{fig:surfaces}). A*
therefore expands the points near this segment first.

\begin{figure}[htbp]
  \centering
  \resizebox{\linewidth}{!}{\def\kD{0.45}\def\sgx{1.6}
\newcommand{\surfone}[4]{%
  \begin{tikzpicture}
    \begin{axis}[land3d, land view={-28}{30}{0.52}, colormap name=#1,
                 point meta min=0, point meta max=3.6,
                 xmin=-3.2, xmax=3.2, ymin=-2, ymax=2, zmin=0, zmax=3.6]
      \draw[gridgray, line width=0.4pt] (-3.2,-2,0) -- (3.2,-2,0)
        -- (3.2,2,0) -- (-3.2,2,0) -- cycle;
      \addplot3[tilesurf, domain=-3.2:3.2, y domain=-2:2,
                samples=17, samples y=11] {#2};
      \path (-3.2,-2,3.6) (3.2,-2,3.6) (3.2,2,3.6) (-3.2,2,3.6);
      \node[sball, minimum size=2.6mm] at (-\sgx,0,#3) {};
      \node[gball, minimum size=2.6mm] at (\sgx,0,#4) {};
    \end{axis}
  \end{tikzpicture}}
\begin{tabular}{@{}c@{\hspace{2pt}}c@{\hspace{2pt}}c@{\hspace{2pt}}c@{\hspace{2pt}}c@{}}
  \surfone{hblue}{\kD*veclen(x+\sgx,y)}{0}{\kD*2*\sgx}
  & \raisebox{1.2cm}{\Large $+$}
  & \surfone{horange}{\kD*veclen(x-\sgx,y)}{\kD*2*\sgx}{0}
  & \raisebox{1.2cm}{\Large $=$}
  & \begin{tikzpicture}
      \begin{axis}[land3d, land view={-28}{30}{0.52}, colormap name=hgreen,
                   point meta min=0, point meta max=3.6,
                   xmin=-3.2, xmax=3.2, ymin=-2, ymax=2, zmin=0, zmax=3.6]
        \draw[gridgray, line width=0.4pt] (-3.2,-2,0) -- (3.2,-2,0)
          -- (3.2,2,0) -- (-3.2,2,0) -- cycle;
        \addplot3[tilesurf, domain=-3.2:3.2, y domain=-2:2,
                  samples=17, samples y=11]
          {\kD*(veclen(x+\sgx,y)+veclen(x-\sgx,y))};
        \path (-3.2,-2,3.6) (3.2,-2,3.6) (3.2,2,3.6) (-3.2,2,3.6);
        \draw[cbGreen!60!black, line width=1.8pt]
          (-\sgx,0,{\kD*2*\sgx}) -- (\sgx,0,{\kD*2*\sgx});
        \node[sball, minimum size=2.6mm] at (-\sgx,0,{\kD*2*\sgx}) {};
        \node[gball, minimum size=2.6mm] at (\sgx,0,{\kD*2*\sgx}) {};
      \end{axis}
    \end{tikzpicture}\\
  {\footnotesize\textcolor{rblue}{$g$}: cost so far} &
  & {\footnotesize\textcolor{cbOrange}{$h$}: estimate of what is left} &
  & {\footnotesize\textcolor{cbGreen}{$f = g + h$}}\\
  {\scriptsize\textcolor{neutralg}{Dijkstra: a cone at $s$}} &
  & {\scriptsize\textcolor{neutralg}{greedy: a cone at $t$}} &
  & {\scriptsize\textcolor{neutralg}{A*: a trough from $s$ to $t$}}
\end{tabular}
}
  \caption{The keys as surfaces over a plane: a cone around $s$, a cone around
    $t$, and their sum, a trough whose floor is the segment from $s$ to
    $t$.}\label{fig:surfaces}
\end{figure}

\subsection{Pseudocode}\label{sec:astar-alg}

Algorithm~\ref{alg:astar} is Algorithm~\ref{alg:dijkstra} with the key $g$
replaced by $g + h$ in the two lines marked ``changed''. With the \emph{zero
heuristic}, $h(v) = 0$ for every $v$, it is exactly Dijkstra's algorithm.
Figure~\ref{fig:astar-flow} shows the loop shared by the three algorithms.

\begin{algorithm}[htbp]
  \caption{A* search}\label{alg:astar}
  \begin{algorithmic}[1]
    \Require graph $G = (V,E,w)$ with nonnegative weights, source $s$, destination $t$,
      heuristic $h$
    \Ensure a path from $s$ to $t$, or failure
    \Procedure{A-Star}{$G, s, t, h$}
      \For{each vertex $v$ of $V$}
        \State $g(v) \gets \infty$; \ $\pi(v) \gets \nil$
      \EndFor
      \State $g(s) \gets 0$
      \State $Q \gets$ a priority queue keyed by $g + h$ that holds $s$
        \Comment{changed}
      \State $S \gets \emptyset$
      \While{$Q$ is not empty}
        \State $u \gets \Call{Extract-Min}{Q}$
        \If{$u = t$}
          \State \Return \Call{Reconstruct}{$\pi, t$}
        \EndIf
        \State add $u$ to $S$
        \For{each edge $(u,v)$ such that $v$ is not in $S$}
          \If{$g(u) + w(u,v) < g(v)$}
            \State $g(v) \gets g(u) + w(u,v)$; \ $\pi(v) \gets u$
            \State \Call{Insert-or-Decrease-Key}{$Q, v, g(v) + h(v)$}
              \Comment{changed}
          \EndIf
        \EndFor
      \EndWhile
      \State \Return failure
    \EndProcedure
  \end{algorithmic}
\end{algorithm}

\begin{figure}[htbp]
  \centering
  \tikzset{
  fstep/.style  = {draw=rblue, fill=rbluelt, line width=0.7pt, rounded corners=2pt,
                   text width=6.2cm, align=center, font=\small, inner sep=4pt},
  fterm/.style  = {fstep, fill=rblue, text=white, rounded corners=8pt,
                   text width=2.6cm, font=\small\bfseries},
  fdec/.style   = {draw=alggreedy, fill=alggreedy!12, line width=0.7pt, diamond,
                   aspect=2.4, align=center, font=\small, inner sep=1pt},
  fout/.style   = {fstep, text width=3.2cm},
  farr/.style   = {-{Stealth[length=2mm]}, line width=0.8pt, draw=rbluedk},
  flab/.style   = {font=\scriptsize\bfseries, text=rbluedk, inner sep=2pt},
}
\begin{tikzpicture}[node distance=0.55cm]
  \node[fterm] (start) {Start};
  \node[fstep, below=of start] (init) {$g(s) \gets 0$, every other $g(v) \gets \infty$\\
    $Q \gets \{s\}$, \ $S \gets \emptyset$};
  \node[fdec, below=of init] (empty) {Is $Q$\\empty?};
  \node[fstep, below=of empty] (pop) {$u \gets$ the vertex of $Q$ with the\\
    smallest key};
  \node[fdec, below=of pop] (goal) {$u = t$?};
  \node[fstep, below=of goal] (close) {add $u$ to $S$};
  \node[fstep, below=of close] (relax) {for each edge $(u,v)$ with $v$ not in $S$:\\
    if $g(u) + w(u,v) < g(v)$, then set $g(v)$ and $\pi(v) \gets u$,\\
    and update the key of $v$ in $Q$};
  \node[fout, right=1.4cm of empty] (fail) {return failure};
  \node[fout, fill=algastar!12, draw=algastar, right=1.4cm of goal] (done)
    {\textsc{Reconstruct}$(\pi,t)$\\(Algorithm~\ref{alg:reconstruct})};
  \draw[farr] (start) -- (init);
  \draw[farr] (init) -- (empty);
  \draw[farr] (empty) -- node[flab, right] {no} (pop);
  \draw[farr] (empty) -- node[flab, above] {yes} (fail);
  \draw[farr] (pop) -- (goal);
  \draw[farr] (goal) -- node[flab, above] {yes} (done);
  \draw[farr] (goal) -- node[flab, right] {no} (close);
  \draw[farr] (close) -- (relax);
  \draw[farr] (relax.west) -- ++(-0.7,0) |- (empty.west);
  \node[draw=cbRed, fill=cbRed!6, line width=0.7pt, rounded corners=2pt,
        text width=3.6cm, font=\small, align=left, inner sep=5pt,
        right=1.0cm of relax] (key)
    {\textbf{key of $v$}\\[2pt]
     Dijkstra: \ $g(v)$\\
     greedy: \ $h(v)$\\
     A*: \ $g(v) + h(v)$};
  \draw[cbRed, line width=0.7pt, dashed] (key.west) -- (relax.east);
  \draw[cbRed, line width=0.7pt, dashed] (key.north) |- (pop.east);
\end{tikzpicture}

  \caption{Flow chart of the search loop. The three algorithms differ only in
    the key of a vertex (red box).}\label{fig:astar-flow}
\end{figure}

\subsection{A* on the running example}\label{sec:worked}

Table~\ref{tab:trace-h} and Figure~\ref{fig:story} show A* on the running
example. In step 2, A* expands \cty{N} (key $2 + 5 = 7$) rather than \cty{G}
(key $1 + 8 = 9$), which Dijkstra's algorithm would take first. The same step
lowers $g(\cty{B})$ from 5 to 3 through \cty{N}, repairing the error of greedy
best-first search. In step 5 the destination is extracted with $g = 8$, and
\cty{G}, \cty{K} and \cty{C}, whose keys exceed 8, are never expanded. A* thus
finds the shortest path with five expansions instead of eight.
Sections~\ref{sec:heuristics} and~\ref{sec:correctness} show when this result is
guaranteed, and Section~\ref{sec:efficiency} explains the saving.

\begin{table}[htbp]
  \centering
  \caption{A* on the running example. The key of a vertex is
    $f = g + h$.}\label{tab:trace-h}
  \small
  \begin{tabular}{@{}c>{\raggedright}P{4.4cm}cc>{\raggedright\arraybackslash}P{3.8cm}@{}}
    \toprule
    Step & $Q$ before extraction (vertex: key) & Extracted & $g + h = f$
      & Updates\\
    \midrule
    \inputrows{figures/trace_original.tex}
    \bottomrule
  \end{tabular}
\end{table}

\begin{figure}[htbp]
  \centering
  \resizebox{0.62\textwidth}{!}{\begin{tikzpicture}[x=1cm, y=1cm]
  \fill[black!10] (-0.529,-0.393) -- (7.052,-1.657) -- (8.132,1.475) -- (0.551,2.739) -- cycle;
  \draw[draw=rblue!14, fill=rbluelt!35, line width=0.6pt] (-0.529,-0.243) -- (7.052,-1.507) -- (8.132,1.625) -- (0.551,2.889) -- cycle;
  \coordinate (D) at (0.432,1.253);
  \coordinate (N) at (2.333,1.598);
  \coordinate (B) at (3.802,0.691);
  \coordinate (H) at (5.702,1.037);
  \coordinate (S) at (7.171,0.130);
  \coordinate (G) at (0.950,0.173);
  \coordinate (K) at (4.234,1.944);
  \coordinate (C) at (5.584,-0.599);
  \draw[draw=algastar, line width=2pt] (D) -- (N);
  \draw[draw=rblue!55, line width=1pt] (D) -- (B);
  \draw[draw=algastar, line width=2pt] (N) -- (B);
  \draw[draw=algastar, line width=2pt] (B) -- (H);
  \draw[draw=algastar, line width=2pt] (H) -- (S);
  \draw[draw=rblue!55, line width=1pt] (D) -- (G);
  \draw[draw=rblue!55, line width=1pt] (N) -- (K);
  \draw[draw=rblue!55, line width=1pt] (K) -- (H);
  \draw[draw=rblue!55, line width=1pt] (B) -- (C);
  \draw[draw=rblue!55, line width=1pt] (C) -- (S);
  \node[isoweight] at ($(D)!0.5!(N)$) {2};
  \node[isoweight] at ($(D)!0.5!(B)$) {5};
  \node[isoweight] at ($(N)!0.5!(B)$) {1};
  \node[isoweight] at ($(B)!0.5!(H)$) {2};
  \node[isoweight] at ($(H)!0.5!(S)$) {3};
  \node[isoweight] at ($(D)!0.5!(G)$) {1};
  \node[isoweight] at ($(N)!0.5!(K)$) {3};
  \node[isoweight] at ($(K)!0.5!(H)$) {4};
  \node[isoweight] at ($(B)!0.5!(C)$) {3};
  \node[isoweight] at ($(C)!0.5!(S)$) {10};
  \isodisc{D}{algastar!15}{algastar}{D}{7}
  \isodisc{N}{algastar!15}{algastar}{N}{5}
  \isodisc{B}{algastar!15}{algastar}{B}{4}
  \isodisc{H}{algastar!15}{algastar}{H}{2}
  \isodisc{S}{algastar!15}{cbRed}{S}{0}
  \isodisc{G}{white}{rblue!70}{G}{8}
  \isodisc{K}{white}{rblue!70}{K}{4}
  \isodisc{C}{white}{rblue!70}{C}{7}
\end{tikzpicture}
}
  \caption{A* on the running example. The green path is the path that A*
    returns, and its cities are exactly the cities that A*
    expands.}\label{fig:story}
\end{figure}
